% This document is made by doing:
%   MAKE-MPI-CHAPTER errata-40.tex
%   use ghostview to strip the first two pages off
%   update change data in docs.html
%   put on home page in compressed and uncompressed format

\def\mf#1{{{\sf MPI\_#1}}}
\def\cf#1{{{\tt #1}}}

% For each batch of errata, we surround it with
% \ifbati (ii, etc)
% \fi % bati (ii, etc)

\def\pending{\textbf{Pending Errata:}}

% Is this a draft or an official version?
\newif\ifdraft
\drafttrue
%\draftfalse
%
% Macro to point at the discussion.  This is soley to help in generating the
% HTML online version
\def\webdiscussion#1{\relax}
%
% This errata contains the current list of both accepted and
% pending items, in page order.  Each has a comment (at the top!) that
% gives the state
% The categories are
%   Original (posted on the forum)
%   May-1-01     (first batch)

\chapter*{Errata for \mpiivdoto/}

This document was processed on \today.
\vskip 0.3truein

The known corrections to \mpiivdoto/ are listed in this document.  All
page and line numbers are for the official version of the \mpiivdoto/
document available from the \MPI/ Forum home page at
\url{www.mpi-forum.org}.  Information on reporting mistakes in the
\mpi/ documents is also located on the \MPI/ Forum home page.

\ifdraft
This is a \textbf{DRAFT} and is not official.  It may contain errata items
that are currently under discussion and have not yet been voted in by
the MPI Forum.  Such items are marked as \pending.
\fi

\begin{itemize}

% Example of an errate item
%\item Page 53, lines 24--25 reads:
%
%\begin{quote}
%Any other type of request is and the request handle is set to \const{MPI\_REQUEST\_NULL}. \mpifunc{MPI\_WAIT} is a non-local operation.
%\end{quote}
%
%\noindent but should be:
%
%\begin{quote}
%Any other type of request is deallocated and the request handle is set to \const{MPI\_REQUEST\_NULL}. \mpifunc{MPI\_WAIT} is a non-local operation.
%\end{quote}
%
\item Page 113, lines 6--7 read:
\begin{verbatim}
     MPI_Psend_init(message, partitions, count, xfer_type, dest, tag,
          info, MPI_COMM_WORLD, &request);
\end{verbatim}
but should be
\begin{verbatim}
     MPI_Psend_init(message, partitions, count, xfer_type, dest, tag,
                    MPI_COMM_WORLD, info, &request);
\end{verbatim}

\item Page 113, lines 25--26 read:
\begin{verbatim}
     MPI_Precv_init(message, partitions, count, xfer_type, source, tag,
           info, MPI_COMM_WORLD, &request);
\end{verbatim}
but should be
\begin{verbatim}
     MPI_Precv_init(message, partitions, count, xfer_type, source, tag,
                    MPI_COMM_WORLD, info, &request);
\end{verbatim}

\item Page 114, lines 28--29 read:
\begin{verbatim}
     MPI_Psend_init(message, send_partitions, count, send_type, dest, tag,
               info, MPI_COMM_WORLD, request);
\end{verbatim}
but should be
\begin{verbatim}
     MPI_Psend_init(message, send_partitions, count, send_type, dest, tag,
                    MPI_COMM_WORLD, info, &request);
\end{verbatim}

\item Page 115, lines 12--13 read:
\begin{verbatim}
     MPI_Precv_init(message, recv_partitions, recv_partlength, MPI_DOUBLE,
               source, tag, info, MPI_COMM_WORLD, &request);
\end{verbatim}
but should be
\begin{verbatim}
     MPI_Precv_init(message, recv_partitions, recv_partlength, MPI_DOUBLE,
                    source, tag, MPI_COMM_WORLD, info, &request);
\end{verbatim}

\item Page 115ff, Example 4.4 has several corrections, including
                    a \texttt{firstprivate(send\_partitions)} at page
                    116, line 24;
                    a \texttt{firstprivate(recv\_partitions)} at page
                    116, line 43, and corrected code for page 117,
                    lines 1--28.
% Fixd with PR 726
\item Page 116, lines 21--22 read:
\begin{verbatim}
     MPI_Psend_init(message, send_partitions, 1, send_type, dest, tag,
               info, MPI_COMM_WORLD, &request);
\end{verbatim}
but should be
\begin{verbatim}
     MPI_Psend_init(message, send_partitions, 1, send_type, dest, tag,
                    MPI_COMM_WORLD, info, &request);
\end{verbatim}

\item Page 116, lines 40--41 read:
\begin{verbatim}
     MPI_Precv_init(message, recv_partitions, recv_partlength, MPI_DOUBLE,
               source, tag, info, MPI_COMM_WORLD, &request);
\end{verbatim}
but should be
\begin{verbatim}
     MPI_Precv_init(message, recv_partitions, recv_partlength,
                    MPI_DOUBLE, source, tag, MPI_COMM_WORLD, info,
                    &request);
\end{verbatim}

\item Page 383, line 36, add\\

Passing \const{MPI\_COMM\_NULL} as \mpiarg{comm} will return the string
\mpistring{MPI\_COMM\_NULL}.\\

Similarly, on page 385, after line 24, add\\

Passing \const{MPI\_DATATYPE\_NULL} as \mpiarg{datatype} will return the string
\mpistring{MPI\_DATATYPE\_NULL}.\\

and similarly, on page 386, after line 21, add\\

Passing \const{MPI\_WIN\_NULL} as \mpiarg{win} will
return the string \mpistring{MPI\_WIN\_NULL}.

%% Errata from PR652

\item Page 407, lines 7--10 reads:\\
\funcargalone{\OUT}{coords}{integer array (of size \mpiarg{maxdims}) containing the Cartesian coordinates of specified process (array of integers)}

but should read\\
\funcargalone{\OUT}{coords}{coordinates of the \MPI/ process with the rank \mpiarg{rank} in Cartesian structure}

The behavior when \mpiarg{maxdims} is less than the number of dimensions of the
Cartesian topology is clarified.

%% Errata from PR 690

\item Page 484, line 23 reads\\
\mpibind{MPI\_Info\_create\_env(\gb{}int~argc, char~argv[], MPI\_Info~*info)}
but should read\\
\mpibind{MPI\_Info\_create\_env(\gb{}int~argc, char~*argv[], MPI\_Info~*info)}

\item Page 521, line 7. Add
\begin{quote}
The \mpi/ procedures described in this section require the \wpm{}, meaning that \mpifunc{MPI\_INIT} or \mpifunc{MPI\_INIT\_THREAD}  has been used to initialize \mpi/.
\end{quote}
to the end of the paragraph in Section 11.7.

\item Page 601, lines 28--31: the definition of \mpifunc{MPI\_WIN\_TEST} was clarified.
The new text reads:

\mpifunc{MPI\_WIN\_TEST} is a local procedure.
Repeated calls to \mpifunc{MPI\_WIN\_TEST} with the same \mpiarg{win} argument
will eventually return \mpiarg{flag}\code{ = true} once all accesses to the
local window by the group to which it was exposed by the corresponding
\mpifunc{MPI\_WIN\_POST} call have been completed as indicated by matching
\mpifunc{MPI\_WIN\_COMPLETE} calls, and \mpiarg{flag}\code{ = false} otherwise.

%% PR 729. This was voted in as an errato, so is added here.

\item Page 607, replace lines 40--44 with:
\begin{quote}
For windows in the separate memory model,
a call to \mpifunc{MPI\_WIN\_SYNC} synchronizes the private and public window
copies of \mpiarg{win} at the calling \MPI/ process, as described in
Section~12.7.

In the unified memory model, \mpifunc{MPI\_WIN\_SYNC} may be used to
order load and store accesses to shared memory and to ensure
visibility of store updates in shared memory
for other threads and \MPI/ processes.

A call to \mpifunc{MPI\_WIN\_SYNC} does not open or close an epoch and
does not complete any pending \RMA/ operations. A call
to \mpifunc{MPI\_WIN\_SYNC} does not
guarantee \mpitermni{progress}\mpitermindex{progress!one-sided
communication}\mpitermindex{progress!load/store access
synchronization} of any pending \MPI/ operation.

\end{quote}

%% PR 747.

\item Page 640, after line 36, is missing \mpifunc{MPI\_Status\_set\_elements\_c}:\\
\mpibind{MPI\_Status\_set\_elements\_c(\gb{}MPI\_Status~*status, MPI\_Datatype~datatype, MPI\_Count~count)}

Several lists of routines need to have this routine added as well.

\item Page 733, line 6.  Add after the end:

\begin{quote}
It is erroneous to
bind a control variable, performance variable, or event to a handle that
would not be valid to use as an input argument to another \MPI/ call (excluding
calls to the \MPI/ Tool Information Interface) at the same point of
execution.
\end{quote}

%% PR 161

\item Pages 743--754. Several functions marked handles as \IN{} rather
than \INOUT{}. These are:

\mpifunc{MPI\_T\_CVAR\_WRITE}, page 743, line 38\\
\mpifunc{MPI\_T\_PVAR\_HANDLE\_ALLOC}, page 751, line 2\\
\mpifunc{MPI\_T\_PVAR\_HANDLE\_FREE}, page 751, line 44\\
\mpifunc{MPI\_T\_PVAR\_START}, page 752, line 17\\
\mpifunc{MPI\_T\_PVAR\_STOP}, page 752, line 36\\
\mpifunc{MPI\_T\_PVAR\_WRITE}, page 753, line 32, and\\
\mpifunc{MPI\_T\_PVAR\_RESET}, page 754, line 8.

%% PR 754, issue 443

\end{itemize}
